%! Boxplots, Outliers, and Resistance — Unit 2, Lesson 2.3 — Student Packet Cover Sheet
% PILOT SHAPE, Unit 2 timing (5 / 45 / 5 / 5). THREE packet rows — Warm-Up, Guided Notes &
% Practice, Homework. The homework is the individual practice, assigned at Close & Assign and
% done outside class, and it is SCORED, so its cell is a \blank{}, never NA.
\documentclass[12pt]{article}
\usepackage{saar-article}
\usepackage{saar-boxes}
\usepackage{ltablex}
\keepXColumns
% Measured banner: the band grows to fit the title block, so a title long enough to wrap grows
% the band rather than running out of it (the stats course's \coverbanner; saar-article does not
% carry one, so it is defined here — every cover copies this block verbatim).
\newsavebox{\bannerbox}
\newlength{\bannerht}
\newlength{\coverbannerpad}
\setlength{\coverbannerpad}{0.16in}       % breathing room above and below the text
\newcommand{\coverbanner}[2]{%
  \sbox{\bannerbox}{%
    \begin{minipage}{\dimexpr\paperwidth-1.4in\relax}%
      \centering
      {\color{white}\textbf{\LARGE \CourseName}}\\[4pt]
      {\color{frostmid}\large\textbf{#1}}\\[6pt]
      {\color{white}\textbf{\large #2}}%
    \end{minipage}}%
  \setlength{\bannerht}{\dimexpr\ht\bannerbox+\dp\bannerbox+2\coverbannerpad\relax}%
  \noindent\begin{tikzpicture}[remember picture, overlay]
    \fill[cerulean] (current page.north west)
      rectangle ([yshift=-\bannerht]current page.north east);
    \node[anchor=north, inner sep=0pt]
      at ([yshift=-\coverbannerpad]current page.north) {\usebox{\bannerbox}};
  \end{tikzpicture}\par
  % The text body starts 0.6in below the page edge (geometry top=0.6in), so reserve only what
  % the band overruns that, plus a gap under the band.
  \vspace*{\dimexpr\bannerht-0.6in+0.16in\relax}%
}

\begin{document}

% Full-bleed cerulean banner, sized to the title block.
\coverbanner{Unit 2}{Lesson 2.3 \quad Boxplots, Outliers, and Resistance}

\vspace{0.06in}
% NAMESTRIP: this is the ONLY component that carries the name row.
\namedateperiod
\vspace{0.18in}

% GRADUAL RELEASE: the vocabulary is taught in the guided notes, so the targets name it
% plainly. The standard codes (PS.DS.1a, PS.DS.2c-d) stay in the teacher-facing plan.
\begin{learningtargetbox}
\small
\begin{enumerate}[leftmargin=1.5em, itemsep=5pt]
  \item \ldots{} read a \textbf{boxplot} and finish one from a \textbf{five-number summary},
        knowing that each of its four sections holds about $25\%$ of the data.
  \item \ldots{} use the \textbf{$1.5 \times \text{IQR}$ rule} to find the \textbf{fences},
        name every \textbf{outlier}, and read a \textbf{modified boxplot}.
  \item \ldots{} read the median, the quartiles, and any \textbf{percentile} off a
        \textbf{cumulative relative frequency graph}, and say what it means in context.
  \item \ldots{} explain why the median and IQR are \textbf{resistant} and the mean and standard
        deviation are not --- and choose which center and spread to report.
\end{enumerate}
\end{learningtargetbox}

\vspace{0.14in}

\begin{tocbox}
\small
\renewcommand{\arraystretch}{1.5}
\begin{tabularx}{\linewidth}{@{} c l X r @{}}
\rowcolor{frost}
\textbf{\#} & \textbf{Component} & \textbf{Description} & \textbf{Score} \\
\midrule
1 & Warm-Up & Pine Street Pizza's Friday deliveries: the median, the quartiles, the IQR, and
             counts turned into percents & \blank{1.2cm} \\
2 & Guided Notes \& Practice & Boxplots, the $1.5 \times \text{IQR}$ rule and outliers, the
             cumulative relative frequency graph, and what one flat tire does to each summary
             --- then one shelter we work together & \blank{1.2cm} \\
% Row 3 is the lesson's GRADED practice, assigned at the close and done outside class. Packet
% pages are the default; a Desmos activity may replace them on a given day, decided at Close &
% Assign. The row and its score cell never change.
3 & Homework & Elm Park Bike Shop: the five-number summary, the fences, a cumulative graph,
             and what selling one \$900 bike does to each summary ---
             \textbf{scored, due the first class after two study halls} & \blank{1.2cm} \\
\midrule
  & \multicolumn{2}{r}{\textbf{Total}} & \blank{1.2cm} \\
\end{tabularx}
\end{tocbox}

\vspace{0.14in}

% Keep in Mind carries the lesson's CONTENT — the rule, the test, the trap — never its process.
\begin{remindbox}
\small
A value is an \textbf{outlier} only if it falls outside a fence: below $Q_1 - 1.5 \times
\text{IQR}$ or above $Q_3 + 1.5 \times \text{IQR}$. Check the fence --- never decide because a
point ``looks far.'' The \textbf{median} and \textbf{IQR} are \textbf{resistant}: one extreme
value barely moves them. The \textbf{mean} and \textbf{standard deviation} are not --- every
value is in their arithmetic. \textbf{Watch out:} an outlier does not move every summary
equally. When data are skewed or have outliers, report the median and the IQR.
\end{remindbox}

\end{document}
